\documentclass[11pt,a4paper]{article}
\usepackage[margin=23mm,headheight=15pt,headsep=7mm,footskip=10mm]{geometry}
\usepackage[T1]{fontenc}
\usepackage[utf8]{inputenc}
\usepackage{amsmath,amsthm}
\usepackage{newtxtext,newtxmath}
\usepackage{booktabs,longtable,array,calc}
\usepackage{microtype,xurl,graphicx}
\usepackage[hidelinks]{hyperref}
\usepackage{fancyhdr}
\usepackage{titlesec}
\setlength{\parindent}{0pt}
\setlength{\parskip}{4pt}
\setlength{\emergencystretch}{2em}
\linespread{1.02}
\setcounter{secnumdepth}{0}
\setcounter{tocdepth}{2}
\newcommand{\tightlist}{\setlength{\itemsep}{0pt}\setlength{\parskip}{0pt}}
\providecommand{\pandocbounded}[1]{#1}
\pagestyle{fancy}
\fancyhf{}
\fancyhead[L]{\footnotesize Ptolemaic rigidity and stability}
\fancyhead[R]{\footnotesize Research supplement v0.3}
\fancyfoot[C]{\thepage}
\renewcommand{\headrulewidth}{0.3pt}
\titleformat{\section}{\large\bfseries}{\thesection}{0pt}{}
\titleformat{\subsection}{\normalsize\bfseries}{\thesubsection}{0pt}{}
\titleformat{\subsubsection}{\normalsize\itshape}{\thesubsubsection}{0pt}{}
\hypersetup{pdftitle={Response and further
development},pdfauthor={},pdfsubject={Research proof supplement; unrefereed}}
\begin{document}
\begin{center}
{\LARGE\bfseries Response and further development\par}
\vspace{6pt}
{\large Supplied v0.2 and research supplement v0.3\par}
\vspace{8pt}
{\small Research version 0.3\quad |\quad 27 September 2026\par}
\end{center}
\vspace{4pt}
\textbf{27 September 2026 --- research supplement v0.3}

\subsection{1. What was accepted and
preserved}\label{what-was-accepted-and-preserved}

The input v0.2 was treated as the reviewed source, not overwritten. Its
full six-job replay passed, including the nested companion exact replay.
Forty input files and 130 nested companion files were manifest-verified.
The saved equality-search summaries changed only in elapsed-time
metadata; the trial records were identical. The new receipt and exact
field diff are retained.

The corrections in the supplied review are preserved: Ercan receives
attribution for coefficient transport and sign-shift; the six-point
extension retains its smaller-cardinality strictness premise; Li--Weston
strict descent is cited where used; the restricted negative-type gap
does not include the automatic centring zero. The v0.2 seven-point
remainder and radius remain the preferred constants for that particular
cardinality.

\subsection{2. New mathematical content}\label{new-mathematical-content}

\textbf{General-cardinality sharp tangent constants.} The seven-point
active-constraint calculation extends to every binding split with
m\textgreater=6 and to the five-point split locally. The constants have
six small exceptions and three simple infinite-family formulas. Eighteen
exact rational-function dual families prove all higher dimensions;
thirty-six rational vectors handle the exceptions. Every active cone
direction is integrated into a genuine Ptolemaic curve, so sharpness is
geometric rather than only linear-programming sharpness.

\textbf{Explicit local error and eigendirection bounds.} A Taylor and
scalar Schur-complement calculation turns the exact cone law into a
two-sided bound with an explicit quadratic remainder and radius for each
cardinality. The general remainder is intentionally loose. It is not
advertised as an improvement of v0.2's sharper seven-point constants.

\textbf{The metric's own critical exponent.} A simple-eigenvalue
implicit-function argument gives the first-order change in the supremal
negative-type exponent. The two sharp coefficients follow exactly from
the cone extremal directions. Strict descent excludes a later return to
negative type. This is a fixed-cardinality result with O(h\^{}2)
remainder, not an unproved uniform expansion over growing dimensions.

\textbf{Small cardinalities and quadratic behaviour.} Three-point
equality is collinearity. Four-point equality is exactly CS(1,3) or
CS(2,2), using the external four-point endpoint, strict descent, Ercan's
sharp scalar estimate and his transport. An unequal-arm perturbation at
the four-point star has gap p(2p-3)t\textsuperscript{2/8+O(t}4), which
rules out a positive linear lower stability constant there. The
five-point equality family is not classified by this work, even though
the local theorem applies to CS(2,3).

\textbf{Collision strata.} The quotient of a colliding pseudometric
identifies the entire kernel at the universal exponent: one zero-sum
condition within each collision class. A separate Gaussian-kernel
estimate gives an explicit subcritical lower gap for a separated smaller
metric. This does not complete a numerical global C\_eta for the full
separated seven-point class.

\subsection{3. Checks that support the new
proof}\label{checks-that-support-the-new-proof}

\begin{itemize}
\tightlist
\item
  Direct symbolic orbit counting verifies the eighteen identities and
  positive coefficient expansions for all k in each half-line. No
  fitting or solver call occurs in the proof verifier.
\item
  A separate Fraction reconstruction checks all small duals and original
  labelled rows, the primal directions and exact feasible curves.
\item
  A separate symbolic calculation reconstructs the four-point
  second-order eigenvalue coefficient from the six distances.
\item
  All 182 seeded perturbation records and eight 70-digit checks are
  retained. Numerical evidence is explicitly corroborative.
\end{itemize}

The certificates are supplied as complete JSON objects and printable
tables, not merely as a PASS statement. The numerical discovery history
remains labelled and separate.

\subsection{4. Remaining review
obligations}\label{remaining-review-obligations}

The local tangent and spectral proofs no longer depend on the companion
all-size theorem, but their interpretation as a classification of the
global equality set still does. Independent specialist review should
focus on the orbit-counting dual identity, its universal positivity
domains, cone realisability and the critical-exponent identification.
Small-cardinality historical priority has not been settled. No
formalisation, public release identifier, public licence or
authenticated external endorsement has been assigned.
\end{document}
